\section{Related Work}
\label{sec:related}

\subsection{Formulaic Alpha Mining with Language Models}

Formulaic alpha mining searches for executable expressions that transform market observations into predictive cross-sectional signals. Published formula libraries demonstrate the interpretability and executability of this representation \citep{Kakushadze2016Formulaic}. Automated approaches then use evolutionary search or learned generators to explore expression spaces and combine the resulting factors \citep{Zhang2020AutoAlpha,Shi2024AlphaForge}. These explicit search spaces support reproducible execution, but they are combinatorial, highly exposed to noisy backtest feedback, and often optimized around a narrow factor statistic rather than the behavior of the final portfolio.

LLMs extend this paradigm by connecting economic hypotheses, symbolic construction, and empirical diagnosis. LLMFactor uses sequential knowledge-guided prompting to extract interpretable factors from financial information \citep{Wang2024LLMFactor}; Automate Strategy Finding with LLM combines multimodal market inputs with a multi-agent portfolio workflow \citep{Kou2024AutomateStrategy}; Alpha-GPT introduces interactive LLM-based alpha generation \citep{Wang2023AlphaGPT}; AlphaAgent regularizes generation for originality, hypothesis--factor alignment, and structural simplicity \citep{Tang2025AlphaAgent}; and QuantaAlpha treats a complete mining trajectory as the unit of evolution and reuses effective sub-trajectories through mutation and crossover \citep{Han2026QuantaAlpha}. These systems establish the value of structured, multi-stage alpha discovery beyond one-shot formula generation. \BenchmarkName\ provides a common contract for evaluating such research procedures as model--harness participants, with executable strategy outputs, explicit resource use, and a frozen final-assessment boundary.

\subsection{Self-Improving Agents and Recursive Search}

Memory and iterative feedback are central to self-improving LLM agents. General agent frameworks such as ReAct interleave reasoning with environment actions \citep{Yao2023ReAct}, while Reflexion and Self-Refine use textual feedback and iterative revision without changing model parameters \citep{Shinn2023Reflexion,Madaan2023SelfRefine}. Voyager demonstrates how an executable skill library can support open-ended, lifelong improvement \citep{Wang2023Voyager}. In finance, FinMem organizes market information with layered memory for trading decisions \citep{Yu2023FinMem}, QuantAgent couples an inner response-refinement loop with an outer knowledge-base update loop \citep{Wang2024QuantAgent}, FactorMiner retains successful patterns, failure constraints, and reusable skills \citep{Wang2026FactorMiner}, and QuantaAlpha preserves trajectory-level evidence for evolutionary refinement \citep{Han2026QuantaAlpha}. AQuA closes bounded loops in which validated experimental evidence guides subsequent proposals, whereas QuantEvolver converts quantitative evaluation into policy-update signals through reinforcement fine-tuning \citep{Guo2026AQuA,Zhang2026QuantEvolver}. Together, these approaches show that improvement can occur through memory retrieval, structured candidate evolution, recursive experimentation, or model-parameter updates.

\BenchmarkName\ treats these forms of adaptation as participant properties. Within-run memory and search policy belong to the registered harness; across-run learning additionally requires a declared development history and a frozen evaluation variant. The benchmark fixes the data boundary, resource definitions, feedback, and financial evaluator so that different adaptive procedures can be compared under stated conditions.

\subsection{Quantitative Benchmarks and Evaluation Protocols}

Existing financial infrastructure serves several different evaluation goals. Qlib standardizes AI-oriented quantitative workflows, while FinRL provides reinforcement-learning environments and trading evaluations \citep{Yang2020Qlib,Liu2020FinRL}. Financial language benchmarks such as FinBen and Finance Agent Benchmark evaluate domain understanding, generation, forecasting, and research tasks \citep{Xie2024FinBen,Bigeard2025FinanceAgent}. General agent benchmarks such as AgentBench and MLAgentBench examine multi-turn interaction, tool use, and iterative experimentation \citep{Liu2023AgentBench,Huang2023MLAgentBench}. AlphaBench focuses on formulaic alpha generation, evaluation, and search \citep{Luo2026AlphaBench}, and recent trading-agent evaluation also examines leakage, attribution, and memory effects \citep{Zhu2026KTDFin}.

AlphaForgeBench is a close comparison: it treats LLMs as quantitative researchers that generate executable factors and strategies, separating model reasoning from backtest execution \citep{Zhang2026AlphaForgeBench}. Its published evaluation uses standardized prompts, independent code generations, and a common backtest engine. Executable outputs and the separation of reasoning from execution are therefore shared design choices. Table~\ref{tab:benchmark_scope} contrasts the documented comparison procedures.

\begin{table}[t]
    \centering
    \footnotesize
    \caption{Protocol scope relative to AlphaForgeBench \citep{Zhang2026AlphaForgeBench}, based on its published pipeline and experimental settings (Sections 3.2 and E.3, version 2). ``Not specified'' describes these passages and does not establish that a capability is absent from its software.}
    \label{tab:benchmark_scope}
    \begin{tabular}{@{}>{\raggedright\arraybackslash}p{0.18\linewidth}>{\raggedright\arraybackslash}p{0.34\linewidth}>{\raggedright\arraybackslash}p{0.40\linewidth}@{}}
        \toprule
        Dimension & AlphaForgeBench & \BenchmarkName \\
        \midrule
        Comparison unit & Independent generations per query and model & Model--harness episodes with prescribed replicates \\
        Research procedure & Standardized prompts and generation settings & Registered harnesses with explicit search and selection policies \\
        Resource boundary & Fixed generation count; cumulative search ledger not specified & Per-replicate admission, reservations, charges, and observed usage \\
        Feedback and selection & Generate-and-backtest pipeline & Authorized research feedback followed by one frozen final strategy \\
        Final assessment & Common historical backtest & Separately authorized sealed final interval \\
        Failure and recovery & API failure disclosed; operation recovery and unknown usage not specified & Failed attempts retained; unresolved operations and measurements explicit \\
        Financial basis & Single-asset strategy backtests across assets & Cross-sectional research cohorts and separate funded accounts \\
        \bottomrule
    \end{tabular}
\end{table}


Our focus is the resource-bounded process that selects a final strategy. Repeated backtest feedback introduces selection bias \citep{Bailey2017Backtest}, while interrupted execution can change the effective search budget if retries generate unrecorded candidates. \BenchmarkName\ makes the model--harness pair, feedback boundary, final commitment, failed attempts, and unresolved measurements part of the evaluation record. Arena recovery preserves operation identity and accepted progress under those rules. These mechanisms define the intended comparison scope; their effect on rankings, completion, and cost requires campaign evidence.
